\documentclass[11pt]{amsart}

\usepackage[T1]{fontenc}
\usepackage{lmodern}
\usepackage{microtype}
\usepackage{amsmath,amssymb,amsthm}
\usepackage[hidelinks]{hyperref}

\hypersetup{
  pdftitle={Counterexamples to an Associated-Prime Conjecture of Fouli, Ha, and Morey}
}

\newtheorem{theorem}{Theorem}
\theoremstyle{remark}
\newtheorem{remark}[theorem]{Remark}

\DeclareMathOperator{\Ass}{Ass}
\DeclareMathOperator{\Min}{Min}

\title[Counterexamples to an associated-prime conjecture]
{Counterexamples to an Associated-Prime Conjecture of Fouli, H\`a, and Morey}

\date{}

\keywords{edge ideal, associated prime, separated stars, cycle}

\begin{document}

\begin{abstract}
We refute the only-if implication in Conjecture~3.7 of Fouli,
H\`a, and Morey on associated primes obtained after adjoining
separated stars to an edge ideal. A single star on the four-cycle
already gives a counterexample, for which we determine the complete
set of associated primes. More generally, over every field and for
every $q\geq 1$, the cycle $C_{3q+1}$ admits $q$ pairwise separated
stars such that the homogeneous maximal ideal is an embedded
associated prime after adjoining the stars, but is not obtained from
any minimal prime of the edge ideal in the conjectured manner.
\end{abstract}

\maketitle

\section{A four-vertex counterexample}

Let $I(G)$ denote the edge ideal of a finite simple graph $G$.
Fouli, H\`a, and Morey conjectured that if $I=I(G)$ and
$f_1,\ldots,f_q$ are separated stars on $I$, then
\begin{equation}\label{eq:FHM}
 P\in\Ass\bigl(R/(I,f_1,\ldots,f_q)\bigr)
 \quad\Longleftrightarrow\quad
 P=(Q,f_1,\ldots,f_q)
\end{equation}
for some $Q\in\Ass(R/I)=\Min(R/I)$; see
\cite[Definitions~2.1--2.2 and Conjecture~3.7]{FHM26}.
The following example disproves the only-if implication in
\eqref{eq:FHM} with one star.

\begin{theorem}\label{thm:C4}
Let $K$ be any field and set
\[
 R=K[x_1,x_2,x_3,x_4],\qquad
 I=(x_1x_2,x_1x_3,x_2x_4,x_3x_4).
\]
Let
\[
 f=x_1+x_2+x_3,\qquad J=(I,f),\qquad
 \mathfrak m=(x_1,x_2,x_3,x_4).
\]
Then $f$ is a star on $I$, centered at $x_1$, and
\[
 (J:x_2+x_3)=\mathfrak m.
\]
Moreover,
\[
 \Ass(R/J)=
 \bigl\{
 (x_1,x_2,x_3),\
 (x_1,x_4,x_2+x_3),\
 \mathfrak m
 \bigr\}.
\]
The first two primes are of the form predicted by
\eqref{eq:FHM}, whereas the embedded prime $\mathfrak m$ is not.
Consequently, the only-if implication in Conjecture~3.7 of
Fouli, H\`a, and Morey is false.
\end{theorem}

\begin{proof}
The generators of $I$ divisible by $x_1$ are $x_1x_2$ and
$x_1x_3$, so $f$ is a star centered at $x_1$. Separation is
vacuous for a single star.

Eliminate $x_1$ using $f=0$ and put
\[
 u=x_2+x_3,\qquad v=x_3,\qquad w=x_4.
\]
Then
\[
 R/J\cong K[u,v,w]/L,
 \qquad L=(u^2,uv,uw,vw).
\]
Since $u\notin L$ and
\[
 (L:u)=(u,v,w),
\]
the class of $u$ has annihilator $(u,v,w)$. Because $f\in J$,
the preimage of this colon ideal in $R$ is $(J:x_2+x_3)$. Hence
$(J:x_2+x_3)=\mathfrak m$ and
$\mathfrak m\in\Ass(R/J)$.

The irredundant primary decomposition
\[
 L=(u,v)\cap(u,w)\cap(u^2,v,w)
\]
yields
\[
 \Ass\bigl(K[u,v,w]/L\bigr)
 =\{(u,v),(u,w),(u,v,w)\}.
\]
Pulling these primes back to $R$ gives the displayed set.
Finally,
\[
 \Min(R/I)=\{(x_1,x_4),(x_2,x_3)\},
\]
and adjoining $f$ to these two primes gives, respectively,
$(x_1,x_4,x_2+x_3)$ and $(x_1,x_2,x_3)$. Thus
$\mathfrak m$ is the additional embedded associated prime.
\end{proof}

\section{Counterexamples with arbitrarily many stars}

\begin{theorem}\label{thm:family}
Fix $n\geq 1$ and a field $K$. Let $G_n=C_{3n+1}$ have vertices
\[
 a_1,c_1,b_1,\ldots,a_n,c_n,b_n,z
\]
in cyclic order
\[
 a_1-c_1-b_1-a_2-c_2-b_2-\cdots-a_n-c_n-b_n-z-a_1.
\]
Set
\[
 R=K[a_1,c_1,b_1,\ldots,a_n,c_n,b_n,z],
 \qquad I=I(G_n),
\]
and let
\[
 f_i=a_i+c_i+b_i\quad(1\leq i\leq n),
 \qquad J=(I,f_1,\ldots,f_n).
\]
Then $f_1,\ldots,f_n$ are pairwise separated stars on $I$. If
\[
 U=\prod_{i=1}^n(a_i+b_i),
 \qquad
 \mathfrak m=(a_1,c_1,b_1,\ldots,a_n,c_n,b_n,z),
\]
then
\[
 (J:U)=\mathfrak m.
\]
Consequently, $\mathfrak m$ is an embedded associated prime of
$R/J$, but
\[
 \mathfrak m\neq(Q,f_1,\ldots,f_n)
 \qquad\text{for every }Q\in\Ass(R/I)=\Min(R/I).
\]
\end{theorem}

\begin{proof}
The support of $f_i$ is the closed neighborhood of $c_i$.
Thus $f_i$ is a star centered at $c_i$. The supports are pairwise
disjoint, and no center is adjacent to a noncenter variable in
another support, so the stars are pairwise separated.

Write $u_i=a_i+b_i$, so $U=u_1\cdots u_n$. For each $i$,
\[
 a_i u_i=a_i f_i-a_i c_i\in J,\qquad
 b_i u_i=b_i f_i-b_i c_i\in J,
\]
and
\[
 c_i u_i=c_i a_i+c_i b_i\in I.
\]
Multiplying by $\prod_{j\neq i}u_j$ shows that every
$a_i,b_i,c_i$ belongs to $(J:U)$.

We also have $zU\in I$. Indeed, every monomial term of $zU$ has
the form
\[
 M=z\epsilon_1\cdots\epsilon_n,
 \qquad \epsilon_i\in\{a_i,b_i\}.
\]
If $\epsilon_1=a_1$, then $za_1\mid M$. Otherwise
$\epsilon_1=b_1$. Unless some edge $b_i a_{i+1}$ with
$1\leq i<n$ divides $M$, induction gives $\epsilon_i=b_i$ for
every $i$, and then $b_nz\mid M$.
Thus every monomial term of $zU$ lies in the monomial ideal $I$.
It follows that
\[
 \mathfrak m\subseteq(J:U).
\]

To prove that the colon is proper, let
\[
 A=K[t_1,\ldots,t_n]/(t_1^2,\ldots,t_n^2)
\]
and define a $K$-algebra homomorphism $\varphi:R\to A$ by
\[
 \varphi(a_i)=t_i,\qquad
 \varphi(c_i)=-t_i,\qquad
 \varphi(b_i)=0,\qquad
 \varphi(z)=0.
\]
Every edge generator of $I$ and every $f_i$ maps to zero, while
\[
 \varphi(U)=t_1\cdots t_n\neq0,
\]
since the squarefree monomials form a $K$-basis of $A$. Hence
$U\notin J$. Since $(J:U)$ is proper and contains the
maximal ideal $\mathfrak m$, we obtain $(J:U)=\mathfrak m$.
Therefore $\mathfrak m\in\Ass(R/J)$. Moreover, the prime ideal
\[
 H=(a_1,c_1,b_1,\ldots,a_n,c_n,b_n)
\]
satisfies $J\subseteq H\subsetneq\mathfrak m$, so
$\mathfrak m$ is embedded.

It remains to show that $\mathfrak m$ is not of the form in
\eqref{eq:FHM}. Let $Q\in\Min(R/I)$. The variables outside $Q$
form a maximal independent set $S$ of $G_n$, hence a dominating
set. Each closed neighborhood in a cycle has three vertices, so
\[
 3|S|\geq 3n+1,
 \qquad\text{and therefore}\qquad |S|\geq n+1.
\]
Modulo $Q$, the degree-one component of $\mathfrak m/Q$ has
dimension $|S|$, whereas the images of $f_1,\ldots,f_n$ span a
space of dimension at most $n$. Hence
$(Q,f_1,\ldots,f_n)\neq\mathfrak m$. Finally, $I$ is squarefree,
so $\Ass(R/I)=\Min(R/I)$.
\end{proof}

\begin{remark}\label{rem:provenance}
Up to relabeling, the linear forms in Theorem~\ref{thm:family}
were recorded by Tran as an initially regular sequence realizing
the depth of $C_{3n+1}$; see \cite[Remark~3.2]{Tran} and the
general construction in \cite[Theorem~3.11]{FHM20}. The contribution
here is the explicit annihilator $(J:U)=\mathfrak m$ and the
resulting counterexamples to Conjecture~3.7. The reverse implication
in \eqref{eq:FHM} is not addressed.
\end{remark}

\begin{thebibliography}{9}

\bibitem{FHM20}
L. Fouli, T. H. H\`a, and S. Morey,
\emph{Initially regular sequences and depths of ideals},
J. Algebra \textbf{559} (2020), 33--57.

\bibitem{FHM26}
L. Fouli, T. H. H\`a, and S. Morey,
\emph{Regular sequences of linear forms on monomial ideals},
arXiv:2606.15125v1 [math.AC], 2026.

\bibitem{Tran}
L. Tran,
\emph{Initially regular sequences on cycles and depth of unicyclic graphs},
arXiv:2309.07286v1 [math.AC], 2023.

\end{thebibliography}

\end{document}